% Français 9117, batch 2 — Gallica views 21–40.
% Pass: Opus 5 (claude-opus-5), 2026-09-14 — first pass, unchecked against
% the views by a human.
%
% What the twenty views hold. One continuous draft of the essay, twenty
% written pages in a row, recto and verso of ten leaves, with no blank, no
% guard leaf and no target among them. Every view is a single page of the
% microfilm, the neighbouring leaf showing at the edge. The hand is
% Germain's working hand: fast, corrected throughout, with interlinear
% rewrites, struck lines and a few marginal additions; from view 33 on the
% corrections thin out and the page is closer to a clean copy. It is prose;
% there is no mathematics set as formulae (view 39 names « l'équation du
% cercle » and « les sinus et cosinus » as an example, in words).
%
%   v. 21      f. 10   The page opens mid-sentence, continuing view 20 (not
%                      in this batch): an author effaces the trace of his
%                      first attempts; the reader looks for amusement or
%                      instruction; hence the old separation of imagination
%                      and reason. A last paragraph (« Le bienfait de
%                      l'imprimerie… ») is crossed out and rewritten on v. 22.
%   v. 22      f. 10v  The division of labour in early science; printing;
%                      sciences, letters and arts inspired by one sentiment;
%                      announcement of the next chapter. End of chapter 1.
%   v. 23      f. 11   « Ch. 2 » and its title, which is the title the BnF
%                      gives the volume; the origin of literature, the unity
%                      of action as a means of success before it was a
%                      precept. The foot of the page is left blank.
%   v. 24–28   ff. 11v–13v  Man, the model of all beings, personifies nature;
%                      the first literature poetic and one with physics;
%                      spirits, one creator, eternity and the limit of
%                      thought; the eternity of the soul and the idea of
%                      immortality.
%   v. 29–32   ff. 14–15v  Systems of the esprit philosophique; the phrase of
%                      the preface of the Encyclopédie on the universe as
%                      « un fait unique » (quoted on v. 30); man as the
%                      centre and final cause of all things; alchemy and
%                      judicial astrology.
%   v. 33–36   ff. 16–17v  F. 16 carries a roman « II » under the folio.
%                      Our habit of judging the nature of things by whether
%                      we can form an idea of them; the unity of the true,
%                      the good and the beautiful; Kant's question, set out
%                      and tested on a man who knows only human things.
%   v. 37–40   ff. 18–19v  The object of the philosophical doubt; Kant's
%                      question contained in hers (is our logic that of
%                      absolute reason?); the example of the circle; the
%                      child who strikes back at the table; « qui pourroit
%                      douter de l'absolutisme de nos nécessités » — the
%                      sentence runs on past view 40 into the next batch.
%
% Numbers read on the leaves. Two series.
%   The Bibliothèque's foliation, one figure per leaf at the top right of
% each recto, thin and grey: 10 (v. 21), 11 (v. 23), 12 (v. 25), 13
% (v. 27), 14 (v. 29, the 4 written over a first figure, read as 14), 15
% (v. 31), 16 (v. 33), 17 (v. 35), 18 (v. 37), 19 (v. 39). It is given as
% \folio{N} on those rectos; versos carry none and are named in the notes
% « verso du folio N ». Under the 16 of v. 33, in thick black ink, a roman
% « II ».
%   An ink page series at the top left of every page, recto and verso, each
% number crossed through with a stroke: 11 (v. 21), 12, 13, 14, 15, 16, 17,
% 18, 19, 20 (v. 30), then on v. 31 a struck number whose 2 is read and
% whose second figure is lost in the stroke, 22 (v. 32), 23, 24, 25, 26, 27,
% 28 (v. 38), 29 (v. 39, read with some doubt, the first figure blotted),
% 30 (v. 40) — one page to a view, running on from batch 1's 1–10.
%
% Names in the margins, in a different, larger hand, underlined with a
% stroke back to the text: « M. » and a name beginning with D that is not
% read (v. 33), « Pennequin » (v. 39). Whose they are, and what they mark,
% is not settled here.
%
% Notation and conventions. Her spelling kept (« connoissance », « tems »,
% « enfans », « parceque », « entr'eux », « ame »), including one « avait »
% on v. 40 among the « avoit » elsewhere. « à priori », « à posteriori » are
% underlined and set in \emph. Struck words are kept in \struck; where an
% interlinear rewrite replaces them the rewrite follows in the running text;
% where the syntax of a rewrite was left unresolved it is left so and noted.
% Small crosses after a word, with nothing on the page for them to refer to,
% are recorded in the notes and not in the text.
%
% What could not be read: struck words under heavy erasure (listed as \ill
% inside \struck), the name in the margin of v. 33, a few interlinear words
% (v. 21, 22, 29, 39), and the verb before « pas douté » on v. 39.
%
% Nothing here was checked against any published transcription.

\documentclass[11pt,a4paper]{article}
% The preamble shared by every transcription of the mathematicians' manuscripts in Gallica.
%
% It rests on one idea: **uncertainty compounds, it does not resolve.** A
% transcription that smooths over a doubtful word has destroyed the only thing
% it was made for — a reader must be able to tell, at every point, what was on
% the page and what was supplied. The apparatus macros below are therefore the
% critical apparatus, not decoration.
%
% The same commands are recognised by `scripts/render.mjs`, which produces the
% site's reading view, and by `scripts/tei.mjs`. A macro added here must be
% added there too, or rendering fails loudly — which is the wanted behaviour.
%
% Comments here are English, like the rest of the repository. The documents
% this preamble sets are French, and that is deliberate: see the skills.

\ProvidesPackage{germain}[2026/09/15 Transcriptions of mathematicians' manuscripts in Gallica]

\usepackage{amsmath,amssymb,amsthm}
\usepackage{mathrsfs}
% Kept from the parent preamble so the renderer's diagram support stays
% available; these manuscripts draw no commutative diagrams, and nothing here uses it.
\usepackage{tikz-cd}
\usepackage[english,french]{babel}
\usepackage{fontspec}

% --- Greek ------------------------------------------------------------------
% The text face stays fontspec's default, Latin Modern, which has no Greek:
% XeTeX drops every glyph it lacks with a line in the log and nothing on the
% page, and the Greek words the manuscripts quote vanished from the PDFs. So
% the Greek and Greek Extended blocks get a XeTeX character class of their own,
% and entering or leaving a run of it switches the family to CMU Serif — the
% same Computer Modern design, with polytonic Greek — and back. Only the
% family changes: a Greek word in \textit or \textbf keeps its shape and series.
% The transitions to and from every other class carry the tokens that class
% has towards an ordinary letter, so French spacing before « ; » or inside
% « » still holds after a Greek word. No grouping, so no brace or \endgroup
% can come out unbalanced; maths is left alone, since XeTeX inserts nothing
% there. Kept identical in `exercices.sty`.
\makeatletter
\newfontfamily\germainGreekfont{cmun}[
  Extension=.otf, NFSSFamily=germaingreek,
  UprightFont=*rm, ItalicFont=*ti, SlantedFont=*sl,
  BoldFont=*bx, BoldItalicFont=*bi, BoldSlantedFont=*bl]
\newXeTeXintercharclass\germain@greek
\def\germain@greekrange#1#2{%
  \@tempcnta=#1\relax
  \loop
    \XeTeXcharclass\@tempcnta=\germain@greek
    \ifnum\@tempcnta<#2\advance\@tempcnta\@ne\repeat}
\germain@greekrange{"0370}{"03FF}
\germain@greekrange{"1F00}{"1FFF}
\def\germain@greekprev{\rmdefault}
\def\germain@greekin{%
  \edef\germain@greekprev{\f@family}\fontfamily{germaingreek}\selectfont}
\def\germain@greekout{\fontfamily{\germain@greekprev}\selectfont}
\def\germain@greekjoin{%
  \XeTeXinterchartokenstate=\@ne
  \@tempcnta=\z@
  \loop
    \ifnum\@tempcnta=\germain@greek\else
      \XeTeXinterchartoks\germain@greek\@tempcnta=\expandafter{%
        \expandafter\germain@greekout\the\XeTeXinterchartoks\z@\@tempcnta}%
      \XeTeXinterchartoks\@tempcnta\germain@greek=\expandafter{%
        \the\XeTeXinterchartoks\@tempcnta\z@\germain@greekin}%
    \fi
    \ifnum\@tempcnta<4095\advance\@tempcnta\@ne\repeat}
% babel-french sets its own classes, and saves and restores the interchar
% state, each time a language is selected: join again after it does.
\addto\extrasfrench{\germain@greekjoin}
\addto\extrasenglish{\germain@greekjoin}
\AtBeginDocument{\germain@greekjoin}
\makeatother

\usepackage{geometry}
\geometry{a4paper,margin=25mm,marginparwidth=28mm}
\usepackage{marginnote}
\usepackage{xcolor}
\usepackage[normalem]{ulem}
\setlength{\emergencystretch}{3em}
\usepackage{hyperref}
\hypersetup{colorlinks=true,linkcolor=black,urlcolor=black,pdfborder={0 0 0}}

% --- Batch metadata ---------------------------------------------------------
% Taken as they stand by the rendering script, which shows them at the head of
% the reading view. They fix what the file claims to cover. The macro names are
% the parent project's, so that the scripts stay shared; \folder here means the
% volume's slug — `fr-9115` — and \pages counts Gallica views.

\newcommand{\folder}[1]{\gdef\grothFolder{#1}}
\newcommand{\batch}[1]{\gdef\grothBatch{#1}}
\newcommand{\pages}[2]{\gdef\grothFirst{#1}\gdef\grothLast{#2}}
\newcommand{\foldertitle}[1]{\gdef\grothFoldertitle{#1}}

% The holder's own shelfmark, printed as they write it: « Français 9115 ».
\newcommand{\shelfmark}[1]{\gdef\grothShelfmark{#1}}

% The dating, copied verbatim from the holder's catalogue — never inferred
% here. « XIXe siècle » is the BnF's; a pass that dates a leaf from its content
% says so in a \note{}, as its own inference.
\newcommand{\dating}[1]{\gdef\grothDating{#1}}

% The status notice, printed in the title block and repeated as a diagonal
% watermark on every page of the PDF. The manuscripts are in the public
% domain; what this line declares is not a right but a quality — an
% unverified first pass by a machine. The skills write the line; a file
% without it gets no watermark, so leaving it out is visible in the output.
\newcommand{\watermark}[1]{\gdef\grothWatermark{#1}}

\gdef\grothFolder{?}\gdef\grothBatch{}\gdef\grothFirst{?}\gdef\grothLast{?}%
\gdef\grothFoldertitle{}\gdef\grothShelfmark{}\gdef\grothDating{}\gdef\grothWatermark{}

% --- The résumé -------------------------------------------------------------
\newenvironment{resume}
  {\small\setlength{\parskip}{0.45em}}
  {\par\medskip\hrule\medskip}

% --- Keywords ---------------------------------------------------------------
% In English, closing the modernised reading's résumé: the modern vocabulary
% under which to search for what the leaves construct. The manifest extracts
% them to tag the volume — this line is the single source of the tags.
\newcommand{\keywords}[1]{\par\smallskip\noindent
  {\footnotesize\textsc{Keywords} --- \textit{#1}}\par}

% --- The critical apparatus -------------------------------------------------

% The Gallica view number, in the margin. This is the anchor that lets a
% disagreement over a formula be settled by looking at *one* image rather than
% a volume of 750. The site uses it to turn the facsimile as the transcription
% is scrolled. A view is one scanned image: usually one side of a leaf.
\newcommand{\page}[1]{%
  \marginnote{\footnotesize\color{gray!70}\textsf{v.\,#1}}%
  \label{page:#1}\ignorespaces}

% The BnF's pencilled foliation, where the leaf carries it and a pass can read
% it: « 198r ». This is what the literature cites — « ff. 198r–208v » — and
% the one line that turns a citation into a view. Recorded, never inferred:
% a folio a pass cannot read is not written.
\newcommand{\folio}[1]{%
  \marginnote{\footnotesize\color{gray!50}\textsf{f.\,#1}}\ignorespaces}

% The modernised reading's coarser counterpart to \page{}: which views a
% section is drawn from.
\newcommand{\pagerange}[2]{%
  \marginnote{\footnotesize\color{gray!70}\textsf{v.\,#1--#2}}%
  \ignorespaces}

% Illegible: never guessed.
\newcommand{\ill}{\textcolor{red!60!black}{[\,\ldots\,]}}

% A doubtful reading: the word is given, and flagged as uncertain.
\newcommand{\uncertain}[1]{\uline{#1}}

% An editorial addition — a missing letter, an implied word.
\newcommand{\add}[1]{\textcolor{blue!50!black}{[#1]}}

% Struck out by the writer. What was crossed out is part of the document.
\newcommand{\struck}[1]{\sout{#1}}

% The transcriber's note, on the state of the page or the mathematics.
\newcommand{\note}[1]{\par\smallskip{\small\color{gray!60!black}\textit{[#1]}}\par\smallskip}

% A marginal note of hers, distinct from the body of the text.
\newcommand{\marginal}[1]{\par\smallskip\noindent\hspace{1em}%
  {\small\color{gray!60!black}#1}\par\smallskip}

% --- Title ------------------------------------------------------------------
% The title block is French, like the documents themselves.

\AddToHook{shipout/background}{%
  \ifx\grothWatermark\empty\else
    \begin{tikzpicture}[remember picture, overlay]
      \node[rotate=52, scale=3.4, align=center, text=gray!18,
            font=\sffamily\bfseries]
        at (current page.center) {\grothWatermark};
    \end{tikzpicture}%
  \fi}

\AtBeginDocument{%
  \begin{center}
    {\large\bfseries Manuscrits de mathématiciens (Gallica) ---
      \ifx\grothShelfmark\empty\grothFolder\else\grothShelfmark\fi}\\[2pt]
    {\itshape\grothFoldertitle}\\[4pt]
    {\small vues \grothFirst--\grothLast
      \ifx\grothBatch\empty\else\ (lot \grothBatch)\fi}\\[2pt]
    \ifx\grothDating\empty\else
      {\small\color{gray!60!black}Datation du catalogue : \grothDating}\\[2pt]
    \fi
    \ifx\grothWatermark\empty\else
      {\small\bfseries\color{red!45!gray}\grothWatermark}\\[2pt]
    \fi
    {\footnotesize\color{gray!60!black}Transcription établie d'après les images IIIF de Gallica.
     Source gallica.bnf.fr / Bibliothèque nationale de France.\\
     Les lectures incertaines sont soulignées, les illisibles notées [\,\ldots\,],
     les ajouts entre crochets ; \textsf{v.} numérote les vues de Gallica,
     \textsf{f.} la foliotation au crayon de la BnF.}
  \end{center}
  \vspace{1em}}

\endinput


\folder{fr-9117}
\shelfmark{Français 9117}
\batch{2}
\pages{21}{40}
\dating{1801-1900}
\watermark{Lecture automatique\\non vérifiée}
\foldertitle{« Considérations générales sur l'état des sciences et des lettres aux différentes époques de leur culture, » par Sophie GERMAIN.}

\begin{document}

\page{21}\folio{10r}
connoissance certaine d'une vérité qu'il est à présent en état de
démontrer. Car \struck{l'ouvrage terminé ne laisse \ill{} apercevoir la
trace des essais qui l'ont précédé et ne} bien loin de là chaque
auteur a mis tous ses soins à faire disparoître les traces de ses
premiers essais \struck{il n'a conservé} pour ne conserver que les
formes propres au sujet. \struck{tout été les demeurent sont subsistent,
seules} Le lecteur \struck{emploie} vient ensuite \struck{qui peut} donc
cherche, suivant les dispositions qui lui sont personnelles, soit un
délassement agréable soit une instruction solide. \struck{choisit peut
telles ou telles genre de lecture} Le titre du livre suffit pour qu'il
soit averti \struck{certain} de n'avoir à faire \uncertain{usage}
\struck{employer} que du degré d'attention qu'il veut employer. Il est
\struck{peu donc porté son \ill{}} naturellement porté \struck{à eux
mêmes \ill{}} \struck{doit être tenté de} à croire que les auteurs qui
\struck{\ill{} offrant à son esprit des productions si différentes ont
en les composant agi} ont écrit ou dans \struck{avec} l'abandon d'une
imagination qui erre en liberté, ou avec l'austère méthode d'une
déduction qui ne permet aucun écart\struck{. ainsi se trouve justifiée
les}; de là cette séparation, jadis si respectée \struck{contractée},
entre le domaine de l'imagination et celui de la raison.

\struck{Le bienfait de l'imprimerie qui ne permet plus assure à l'esprit
humain la jouissance \uncertain{de tous} ce que les générations
précédentes ont accumulé d'observations, de comparaisons, de théorie,
de vérités incontestables, augmente incessamment ses forces réelles, il
peut déjà s'élever \struck{de} à la considération des idées générales,
non plus seulement en suivant la pente naturelle d'un goût inné et au
risque de \struck{ne rencon}}
\note{folio 10 de la Bibliothèque, en haut à droite, d'un chiffre fin et gris ;
en haut à gauche, à l'encre, le nombre « 11 » barré d'un trait. Feuillet de brouillon
repris partout. Le haut de la page est très raturé : « Car » est écrit
au-dessus de « l'ouvrage terminé », biffé d'un trait épais ; « bien loin
de là » est écrit sur « et ne » ; « chaque » est porté dans la marge de
gauche en regard de « l'auteur » ; « les traces de ses premiers essais »
est ajouté au-dessus de la fin de la ligne ; « pour ne conserver que »
est porté dans la marge de gauche, et « vient ensuite » et « donc »
au-dessus de la ligne ; la syntaxe de ce passage n'est pas arrêtée
(« Le lecteur vient ensuite donc cherche »). Les mots « suivant les
dispositions qui lui sont personnelles », écrits sur la ligne qui suit
« une instruction solide », sont entourés d'un trait qui les ramène
après « cherche », où ils sont ici mis. « averti », « faire
\uncertain{usage} », « respectée » sont écrits au-dessus des mots biffés
qu'ils remplacent. Dans la marge de gauche, en regard de « offrant à son
esprit », la mention « N. 2 » dans un cadre. Le dernier alinéa, depuis
« Le bienfait de l'imprimerie », est barré d'une grande croix et
s'interrompt sur « au risque de » ; « déjà » est répété dans la marge
de gauche en regard de « il peut déjà ».}

\page{22}
Disons aussi que dans un tems déjà éloigné, l'extrême division du
travail, nécessaire à la science naissante, avoit dû accréditer l'idée
de la spécialité dans les facultés de \struck{l'esprit} l'ame; mais,
aujourd'hui que le bienfait de l'imprimerie, en assurant à l'esprit
humain la jouissance de tout ce que les générations précédentes ont
accumulé, d'observations, de comparaisons, de théories, de vérités
incontestables, il n'aura plus à \struck{refaire} faire les premiers
pas; ses forces réelles augmentent \struck{\ill{}} chaque jour; et
déjà nous nous trouvons ramenés, par la voie sure d'une instruction
approfondie, vers ces idées de simplicité et d'unité qui furent
autrefois des \ill{} révélations du génie, \struck{devinant sa propre
nature.} devinant sa propre nature et cherchant à en étendre les lois
sur l'univers entier.

\struck{Car} N'en doutons plus, \struck{à toutes les époques de leur
développement} les sciences les lettres et les beaux arts ont été
inspirés par un seul et même sentiment. Ils ont reproduit, suivant les
moyens qui constituoient l'essence de chacun d'eux, des copies sans
cesse renouvelées de ce modèle inné, type universel de vérité, si
fortement empreint dans les esprits supérieurs.

Dans le chapitre suivant nous \struck{jetterons} verrons, en jetant un
coup d'œil sur l'histoire de l'esprit humain, \struck{et nous
verrons} comment jusque dans ses écarts mêmes, et en vertu des lois de
son être, tous ses efforts ont été dirigés vers l'ordre la simplicité
et l'unité de conception.
\note{verso du folio 10 ; en haut à gauche, à l'encre, le nombre « 12 » barré d'un trait. Ce verso
récrit l'alinéa barré au bas de la vue 21. « Disons aussi que » est
écrit au-dessus de la première ligne ; « dû » est porté dans la marge de
gauche ; « l'ame » est écrit au-dessus de « l'esprit » biffé ; les mots
écrits au-dessus de « augmentent » et de « des » sont chargés d'encre et
ne se lisent pas ; « verrons, en jetant » est écrit au-dessus de
« jetterons » biffé. Un crochet ouvre, dans la marge de gauche, le
dernier alinéa. La construction « aujourd'hui que \ldots{} il n'aura
plus » est restée telle.}

\page{23}\folio{11r}
\section{Ch.~2. Considérations générales sur l'état des sciences et des lettres aux différentes époques de leur culture.}

\marginal{Variante}
Si, \struck{par} maintenant nous voulons remonter jusqu'à l'origine de
la littérature, \uncertain{nous trouvons qu'elle} a commencé la première
fois que, sortant du cercle étroit des intérêts personnels, l'homme a
essayé de communiquer à ses semblables des \struck{idées qui
n'avoient} sentimens et des idées qui n'avoient aucun but
\struck{utile} utiles.

Le récit des événemens remarquables, la peinture des grandes scènes de
la nature, n'étoient encore que de simples copies de choses
existantes. Lorsque \struck{cessant} au lieu de s'astreindre à
\struck{redire} faire le récit \struck{tel ordre} de certains faits, ou
à dévoiler un certain état \struck{tel ordre} de choses, l'homme de
génie \struck{est} est parvenu à reproduire des impressions reçues
d'ailleurs, à l'aide d'une action dont il avoit imaginé les
\struck{\ill{}} ressorts, \struck{il s'étoit} il s'étoit déjà
\struck{\ill{}} élevé jusqu'à la notion abstraite de l'ordre, pour y
puiser la première des règles de la composition. Il voulut disposer de
\struck{\ill{}} l'attention des autres hommes: \struck{la \ill{}}
l'unité d'action, l'unité d'intérêt, \struck{l'ordre} \struck{\ill{}}
la clarté d'exposition ont été pour lui des moyens de succès
\struck{longtems} avant que l'esprit d'examen en eût fait \struck{ils
ayent été reconnus} comme des préceptes de l'art.
\note{folio 11 en haut à droite ; en haut à gauche, à l'encre, le nombre « 13 » barré d'un trait.
« Ch. 2 » est écrit au-dessus du titre et souligné. « Variante » est
porté dans la marge de gauche, en regard du premier alinéa, suivi d'un
mot biffé. « de » est porté dans la marge devant « La littérature »,
dont la majuscule n'a pas été corrigée ; « nous trouvons qu'elle » est
écrit au-dessus de « sortant du cercle », et la phrase n'a pas été
remise d'accord avec lui. « de certains » et « ou à dévoiler un certain
état » sont écrits au-dessus des mots biffés ; les mots « des impressions
reçues d'ailleurs, » sont entourés d'un trait. « il s'étoit déjà », « l'esprit
d'examen en eût fait » sont écrits au-dessus des lignes. Dans la marge
de gauche, en regard de « l'attention des autres hommes » et des trois
lignes suivantes, une addition de six lignes est entièrement biffée et
ne se lit pas. Le bas de la page est laissé blanc.}

\page{24}
Jeté sur la terre, au milieu de l'immensité des êtres, frappé à la fois
par le spectacle d'une \struck{infinie nombre} infinité de merveilles,
l'homme n'a rien trouvé au dehors de lui de plus merveilleux que lui
même. Il a étendu son existence sur tout ce qui l'environnoit. Son
\struck{\ill{}} individualité lui a été d'abord connue; cherchant
partout sa propre image il a personnifié les êtres inanimés, des êtres
intellectuels enfans de son imagination \struck{ont} \uncertain{ont}
présidé à tous les actes, à tous les phénomènes de \struck{l'ordre
naturel} l'ordre naturel.

Ainsi se manifestoient \struck{déjà} déjà, à cette première époque de
la culture intellectuelle \struck{le sentiment} le sentiment
\struck{commun} profond d'un lien commun entre tous les êtres, et celui
d'un type universel, empreint dans l'intelligence humaine pour lui
servir de \struck{guide} modèle.

Les sciences n'existoient pas encore; mais le besoin d'expliquer
s'étoit fait sentir. La première des littératures fut poétique; ce qui
tenoit lieu des sciences physiques n'étoit pas moins poétique que la
littérature elle même; ou plutôt, ces deux branches du savoir,
\struck{aujourd'hui} tellement séparées aujourd'hui qu'il faut de
l'attention pour remarquer ce qu'elles ont de commun, étoient, dans ces
premiers tems, \struck{la} entièrement confondues.
\note{verso du folio 11 ; en haut à gauche, à l'encre, le nombre « 14 » barré d'un trait. Un crochet ouvre le premier alinéa dans la marge. « Son » est
écrit au-dessus d'un mot biffé devant « individualité » ; le mot écrit
au-dessus de « ont » biffé, devant « présidé », n'est pas sûr. « profond »,
« déjà » et « modèle » sont écrits au-dessus des mots qu'ils remplacent.
Un chiffre « 1 » est porté au-dessus de « êtres inanimés », et de
petites croix suivent « imagination », « sentir », « poétique » et
« premiers tems » : ce sont des signes de renvoi dont l'objet n'est pas
sur cette page.}

\page{25}\folio{12r}
Qu'importoit en effet, à l'égard du caractère de la composition, que
le sujet fût l'homme lui même ou quelqu'un des dieux, demi-dieux, ou
genies qu'il avoit dotés de l'intelligence et des passions humaines?
Des êtres si pareils pouvoient même agir de concert sans nuire à
l'homogénéité d'invention; le merveilleux les unissoit.

Nous apercevons, dans ces premiers essais de la pensée, le goût des
idées générales et le sentiment d'analogie, qui se reproduiront dans la
suite, sous les formes les plus variées. L'individualité et
l'intelligence de l'homme, en vertu de laquelle ses actions sont
dirigées vers le but qu'il veut atteindre lui ont été connues en même
tems que sa propre existence.
\marginal{Dès qu'il porte ses regards autour de lui \uncertain{que}
cherche-t-il? ce qu'il a trouvé en lui même.}
Il voit dans les actes de la nature un ordre et une succession qui lui
paroissent tendre vers un but déterminé; il \struck{attribue} n'imagine
\uncertain{d'autre} cause que l'action d'une intelligence et d'une
volonté; \struck{ces} Cette intelligence, cette volonté, il ne peut les
concevoir sans les attribuer à un être quelconque. Il imagine des êtres
invisibles, parceque en effet il n'en voit aucun. ce sont, suivant
l'importance des actes qu'\struck{il} leur attribue, des dieux, des
demi dieux, ou seulement des genies subalternes. Ces êtres sont amis ou
ennemis; ils combattent entr'eux ou ils unissent leurs forces; ils ont
nos affections, nos peines, nos passions, nos intérêts. Ils sont faits
à notre image. Et pourtant, nous ne pouvons ni les voir, ni les
entendre, ni les palper; ils sont donc immatériels; ce sont des
esprits. \struck{L'homme} fidèle à sa pensée constante, l'homme n'a
jamais cessé de regarder \struck{son} son existence propre comme
\note{folio 12 en haut à droite ; en haut à gauche, à l'encre, le nombre « 15 » barré d'un trait. L'addition marginale « Dès qu'il porte
ses regards\ldots{} en lui même » est appelée par une croix après
« sa propre existence ». Au-dessus de « n'imagine », lui-même écrit
au-dessus d'« attribue » biffé, un mot barré, peut-être « suppose »,
n'est pas sûr ; la première lettre de « d'autre » est empâtée. De petites
croix de renvoi suivent « quelconque », « aucun », « image » et
« propre », sans objet sur la page. La phrase se poursuit sur la vue
suivante.}

\page{26}
comme le type de toutes les autres existences. Après s'être dit: les
esprits existent, ils connoissent, ils veulent, ils agissent
\struck{sur les êtres matériels} et leurs actions \struck{se}
manifestent par les changemens matériels qu'ils opèrent, il
\struck{ne pouvoit tarder} devoit \struck{\ill{} longtems} chercher en
lui même quelque chose de semblable. Nos connoissances, nos volontés et
le principe de nos actions ont donc été attribués à une substance
immatérielle, qui, suivant la diversité \struck{des attributs} de ces
opérations a reçu \struck{des noms} différens noms.

\struck{Nous trouvons, dans} Cette \struck{première} ébauche de nos
connoissances, nous montre l'origine \struck{d'un grand nombre} de la
plupart des idées qui ont été reproduites depuis. La littérature a
conservé les fictions qui furent regardées autrefois comme des
réalités; les sciences physiques ont recueilli les observations que
ces fictions \struck{avoient servi à} expliquoient; la philosophie y a
puisé ses systèmes, et les religions y ont trouvé les élémens de leurs
croyances.

Sans nous astreindre à aucun ordre historique, suivons la marche de
l'esprit humain.

Les observations se sont multipliées; la régularité des mouvemens
célestes et la constance des phénomènes sublunaires ont décélé des lois
immuables. \struck{\ill{}}

Les volontés d'une multitude de personnes n'ont pas ce caractère.
\struck{mais} Un seul homme peut avoir des volontés relatives à
\struck{plusieurs genres d'actions} à des objets différens.
\note{verso du folio 12 ; en haut à gauche, à l'encre, le nombre « 16 » barré d'un trait. « s'être » est écrit au-dessus de « dit » ; les mots « les
esprits existent, ils connoissent, ils veulent, ils agissent » sont
soulignés d'un trait ondulé. « devoit » est écrit au-dessus de mots
biffés, dont le second ne se lit pas ; « de ces opérations », « nous
montre » et « de la plupart » sont écrits au-dessus des lignes ; une
tache d'encre couvre « première ». Au-dessus de « avoient servi à »
biffé, un mot lui-même biffé (peut-être « servoient ») ; « expliquoient »
est écrit par-dessus « expliquer ». « tous » est écrit au-dessus de
« des lois ». Un crochet ouvre dans la marge le dernier alinéa, qui se
termine par un trait.}

\page{27}\folio{13r}
s'il étoit chargé de diriger à la fois plusieurs genres d'actions, il
établiroit un ordre constant qui le dispenseroit d'une attention de
détail. L'état de la société présentoit des \struck{exemples
semblables} exemples. L'homme a dit alors: « un seul être a voulu
l'univers, et il le gouverne; ses volontés sont immuables. »

Nous voyons naître nos semblables; nous avons commencé; l'univers a
donc eu aussi un commencement.

Nous avons une ame immatérielle; elle est la force motrice qui produit
nos actions; l'être des êtres est immatériel; il a créé toutes choses,
et il \struck{les gouverne} agit sur elles.

Il a créé l'univers; il existoit donc avant cet univers. L'esprit
humain étoit arrivé aux limites des analogies, au delà, \struck{il
manquoit de modèle}, il n'avoit plus aucune idée de l'être, car il
manquoit de modèle; cette négation d'idée, cette limite de la pensée a
été exprimée; l'infini est son nom; s'il s'agit de la durée c'est
l'éternité. Le créateur de l'univers \struck{est éternel}
\struck{n'a pas commencé} il n'a pas commencé; il ne doit pas finir;
il est éternel.

Dans ce qui précède, on ne voit pas clairement comment on a été
conduit à \struck{cette} la dernière partie de la proposition: il ne
doit pas finir. \struck{nous n'avons pas suffisamment établi expliqué}

Le voici. Nous \struck{voyons} apercevons au commencement et à la fin
\struck{finir} d'existences pareilles \struck{semblables} à la notre;
ces deux époques sont les limites de la vie: en deçà et au delà nous
trouvons le tems qui ne leur appartient pas; mais nous
\struck{concevons} voyons que d'autres existences \struck{pourront y
trouver place} en remplissent. Ce genre de limites est donc relatif:
\struck{la} ainsi la durée de notre vie est
\note{folio 13 en haut à droite ; en haut à gauche, à l'encre, le nombre « 17 » barré d'un trait, le trait se confondant avec la barre du 7. Au-dessus de « L'état », trois petits
mots écrits entre les lignes ne se lisent pas sûrement ; « état » semble
touché d'un trait. Des guillemets ouvrent « un seul être » et ferment
« immuables ». Deux crochets, dans la marge, ouvrent les alinéas « Nous
voyons naître » et « Nous avons une ame », un troisième « Le voici ».
« car il manquoit de modèle » est écrit au-dessus de « l'être; cette
négation ». « Dans ce qui précède, on ne voit pas clairement comment »
est écrit au-dessus de la ligne, « voit » au-dessus de deux mots biffés, et
« on a été conduit à » est porté dans la marge de gauche. « apercevons
au », « et à la fin », « pareilles », « leur », « voyons », « en » et
« remplissent » sont écrits au-dessus des mots biffés qu'ils
remplacent ; « ainsi » au-dessus de « la » biffé. La phrase se poursuit
sur la vue suivante.}

\page{28}
comprise entre deux limites \struck{absolues} de même genre
\struck{entre deux limites relatives. Conformément aux} L'analogie
vouloit \struck{donc} que l'existence dont on avoit reculé l'origine
jusqu'à la limite absolue fût aussi comprise entre deux limites de même
genre: elle ne devoit pas \struck{avoir de} finir, puisqu'elle n'avoit
pas commencé.

L'intelligence \struck{l'\ill{} humaine} s'étoit déjà approprié la
spiritualité; \struck{elle} elle \struck{ne pouvoit manquer} devoit
prendre aussi possession de l'éternité \struck{son ame auroit dû}
\struck{elle en suivant} sa méthode habituelle, nous l'avons déjà
\struck{vu conformément à}, \struck{suivant laquelle} est
\struck{après avoir} de transporter hors d'elle les lois de sa propre
existence, \struck{elle} de chercher dans les analogies ce qui manque
encore à ses notions intellectuelles et de reporter ensuite vers elle
même les supplémens dont les objets extérieurs lui ont donné l'idée:
\struck{devoit la conduire à l'éternité son ame est elle éternelle}
\marginal{\struck{Cette méthode} Elle se trouvoit donc naturellement
conduite à l'éternité de l'ame}.
On sait qu'en effet cette opinion, présentée sous différentes formes a
eu de nombreux partisans. Cependant \struck{des raisons secondaires ont
fait prévaloir} l'idée moins analogique, de la simple immortalité a
prévalu.

Ajoutons encore quelques observations à ce qui précède.

Lorsque l'individualité multipliée des êtres invisibles satisfaisoit son
imagination, l'homme n'avoit pas encore pratiqué les différens arts en
vertu desquels il assigne aux ouvrages de ses mains une destination
conforme à ses volontés; les procédés mécaniques lui apprirent qu'après
avoir transformé les agens naturels, il pouvoit aussi leur imprimer un
mouvement plus ou moins durable; l'analogie le conduisoit ainsi
\struck{aussi \ill{}} à penser que l'être unique qui gouvernoit le
monde en étoit l'architecte.
\note{verso du folio 13 ; en haut à gauche, à l'encre, le nombre « 18 » barré d'un trait. Un crochet ouvre, dans la marge, l'alinéa
« L'intelligence », et deux autres « Ajoutons encore » et « Lorsque ».
« L'intelligence » est porté dans la marge, devant les mots biffés
qu'il remplace ; « devoit », « de » (trois fois) et « est » sont écrits
au-dessus des lignes ; « nous l'avons déjà » est porté dans la marge de
gauche, suivi de mots biffés, et la syntaxe de la phrase n'est pas tout à
fait arrêtée. L'addition « Elle se trouvoit donc naturellement conduite
à l'éternité de l'ame » est écrite dans la marge de gauche et appelée par
une croix après « lui ont donné l'idée » ; les mots entre les lignes
qui l'accompagnent (« du reste », « de l' ») ne sont pas sûrs. « ainsi »
est écrit au-dessus de « aussi » biffé.}

\page{29}\folio{14r}
Voici une remarque assez singulière: on avoit été mené directement à
dire que le créateur de l'univers n'a pas commencé; l'idée qu'il ne
doit pas finir est pour ainsi dire, symétrique de la première. Eh bien,
en \struck{transportant en lui même} s'appropriant le genre de limite
que son esprit avoit atteint, l'homme ne l'adopte plus pour origine,
mais il en fait le terme de son existence immatérielle. Cette espèce de
paradoxe \struck{se} trouvera son explication lorsque nous nous
occuperons de la liaison établie entre la morale et les croyances.

Nous venons de tracer la plus simple marche que l'intelligence humaine
ait pu suivre. Témoin des merveilles de la nature, voulant, parcequ'elle
en sentoit le besoin, la simplicité l'ordre et les proportions dans ses
propres ouvrages; elle attribue \struck{et} l'unité l'ordre et les
proportions qu'elle remarque dans l'univers à la volonté du créateur.

Mais l'esprit philosophique ne pouvoit se contenter d'une explication
également applicable aux faits les plus contraires. \struck{il est
dans} par son essence \struck{de la volonté d'être} en plus ou moins
arbitraire; il \struck{lui} falloit à l'esprit philosophique un plus
ferme appui; il cherchoit partout les lois de la nécessité: tantôt la
toute puissance divine fut assujettie à de telles lois; tantôt la
matière elle même et ses divers accidens furent regardés comme
nécessaires. \struck{C'est ainsi que l'imagination sans cesse tourmentée
en fait \ill{} essaya \ill{} une infinité d'opinions manières
différentes \ill{} manières diverses de pour se rendre compte des
connoître les réalités dont elle étoit frappée.}

Pressé de trouver au dehors les ressemblances à son modèle intérieur,
\struck{trop} encore peu informé des vérités de la nature; ignorant et
la quantité de chaque phénomène et leurs rapports entr'eux; l'homme de
génie inspiré par le
\note{folio 14 en haut à droite, le 4 écrit par-dessus un premier chiffre ;
en haut à gauche, à l'encre, le nombre « 19 » barré d'un trait. « Voici » est encadré d'un trait. « s'appropriant », « ne », « par son »,
« en », « l'esprit philosophique » (porté dans l'interligne et la marge,
au-dessus de « lui » biffé) et « encore » sont écrits au-dessus des lignes.
La construction « par son essence en plus ou moins arbitraire » est
restée telle. Les quatre lignes biffées de la fin du troisième alinéa
sont reprises plusieurs fois entre les lignes, et plusieurs mots n'en
ont pas été lus. Un trait suit « volonté du créateur ». La phrase se
poursuit sur la vue suivante.}

\page{30}
par le sentiment profond des conditions de l'être, se sentoit entraîné
vers la recherche d'une dépendance mutuelle entre les faits dont il
étoit le témoin. Il les coordonnoit suivant ses convenances
intellectuelles, \struck{il supposoit} \uncertain{et} trouvoit dans les
analogies ce qui manquoit encore à ses connoissances positives.
L'esprit de système \struck{et} devoit égarer l'intelligence
humaine; \struck{par et pourtant il étoit} sans doute \struck{inspiré
par l'intime persuasion} c'étoit \struck{\ill{}} l'effet inévitable de
son penchant à mettre \struck{puisqu'il mettoit} à la place des
certitudes qui n'étoient pas acquises, mille conjectures hardies, qui
se \struck{trouvant} ensuite démenties par des observations nouvelles,
\struck{laissoient l'empreinte des préjugés} léguoient aux générations
suivantes des \struck{idées} véritables préjugés à l'égard des faits
encore inconnus. Il est pourtant certain que cet esprit n'a jamais
cessé d'être guidé par la prévision de la vérité.

Dans ces derniers tems, on a voulu recueillir les différentes branches
de la science en un seul faisceau. L'auteur de la préface de
\struck{cet grand ouvrage} l'Encyclopédie dit, en la terminant:
« l'univers, pour qui sauroit l'embrasser d'un seul coup d'œil,
seroit un fait unique, une grande vérité. »

Ces paroles remarquables renferment le secret des efforts de l'esprit
humain.
\note{verso du folio 14 ; en haut à gauche, à l'encre, le nombre « 20 » barré d'un trait. « dans les analogies ce qui manquoit encore » est
écrit au-dessous de « ses convenances intellectuelles », souligné et
ramené par un trait, et « à ses connoissances positives » est porté dans
la marge de gauche ; « trouvoit » est écrit au bout de la ligne,
au-dessus de mots biffés. « sans doute » est écrit au-dessus de
« pourtant », les mots « c'étoit\ldots{} l'effet inévitable de son
penchant à mettre » au-dessus de la ligne biffée qu'ils remplacent,
« qui n'étoient pas acquises » au-dessus de « certitudes », et
« l'Encyclopédie » au-dessus de « cet grand ouvrage » biffés ; ces
additions et « en un seul faisceau », ajouté sous « recueillir les
différentes branches de la science », sont d'une écriture plus fine et
plus droite que le texte, peut-être d'une autre main. « ensuite » et
« démenties » sont entourés d'un trait. Les guillemets de la citation
sont répétés en tête de chaque ligne ; des traits suivent « en la
terminant » et « une grande vérité ».}

\page{31}\folio{15r}
Chacun des systèmes qu'il a enfantés, avoit pour but de concentrer les
faits alors connus \struck{aux époques qui les ont vu naître} en un
fait unique. On établissoit entre ces faits la relation de cause à
effet; on vouloit une raison pour qu'ils fussent. On cherchoit une
unité, des rapports, un ordre, des proportions, parceque \struck{de
telles} ces conditions sont le caractère du vrai.

\struck{Mais} On n'étoit pas en état de leur donner un appui solide;
mais on généralisoit, avec plus ou moins de bonheur, les résultats dont
on avoit acquis la certitude. Une vérité connue imprimoit alors son
caractère propre à un vaste système; et mille suppositions combloient
ensuite l'intervalle entre cette vérité et celles qui, placées dans un
rang secondaire, sembloient devoir y être liées.

Une idée dominante se retrouve partout; l'homme s'est cru le modèle de
tous les êtres, le but vers lequel ils tendent, le centre de l'univers.
Non seulement ses convenances intellectuelles devoient être réalisées en
toutes choses, mais encore ses moindres convenances usuelles étoient la
cause finale des êtres les plus éloignés de lui. Ainsi le soleil, la
lune, les étoiles sans nombre et presque invisibles, sont faits tout
exprès pour fertiliser ses champs et éclairer ses veilles; le moindre
brin d'herbe développé sans culture, l'animal du désert, le coquillage
qui habite le fond des mers ont leur utilité, ou en d'autres termes, ils
sont faits pour l'homme.

On cherchoit l'unité, l'ordre; on les \uncertain{concentroit} dans des
relations imaginaires. Il y eut d'abord erreur du jugement; mais
l'amour propre enchérissant
\note{folio 15 en haut à droite ; en haut à gauche, à l'encre, un nombre
barré dont on lit le 2 ; le second chiffre se confond avec le trait qui
le barre et n'est pas lu. « alors » et « ces » sont écrits au-dessus des lignes ; des
croix de renvoi suivent « enfantés », « fussent » et « veilles », sans
objet sur la page. Dans « on les \uncertain{concentroit} », la fin du
verbe est surchargée. Un trait suit « faits pour l'homme ». La phrase se
poursuit sur la vue suivante.}

\page{32}
\struck{Cette} erreur, et les religions la consacrèrent bientôt.
\struck{Encore aujourd'hui} Nous voyons encore aujourd'hui la trace des
opinions qui rapportent toute autre existence à celle de l'homme.

Cependant nous appelons présentement fausses ces deux branches
\struck{du} de l'ancien savoir \struck{aucune} \struck{\ill{}} qui,
sous les noms d'alchimie et d'astrologie judiciaire, ont joui pendant
long tems de la plus haute estime.

La première enseignoit que le corps humain est l'abrégé de l'univers.
Les \struck{différentes} diverses substances qu'elle soumettoit à ses
opérations avoient reçu \struck{des} les noms des différens organes
avec lesquels elles avoient des \struck{\ill{}} ressemblances
prétendues. Le foie de souffre est encore connu dans le langage
vulgaire. Cette science vouloit \struck{aussi} aussi l'unité; car elle
cherchoit la panacée ou le remède universel; l'alcali ou le dissolvant
universel; et ce dissolvant devoit finir par réduire toutes autres
substances en un seul élément, qui étoit l'eau. Les métaux \struck{\ill{}}
étoient l'objet de mille doctrines singulières: \struck{ils \ill{}} on
établissoit des rapports entr'eux, et les planètes, dont on leur avoit
donné les noms.

L'astronomie judiciaire apprenoit à juger l'influence des astres sur le
sort de chaque individu. L'homme, persuadé de son importance, se croyoit
menacé par l'apparition des comètes. Les grands de la terre,
renchérissant sur l'amour propre de leurs semblables, \struck{étoient}
ne voyoient \struck{rien du} pas d'évènement plus remarquable que
\struck{leur disparition de} leur propre mort; aussi ne doutoient-ils
pas qu'elle ne fût annoncée par ces astres vagabonds qui bien
certainement
\note{verso du folio 15 ; en haut à gauche, à l'encre, le nombre « 22 » barré d'un trait.
« bientôt » est écrit au-dessus de « Cette » biffé, et « encore
aujourd'hui » au-dessus de « voyons la trace » ; un crochet, devant
« Nous voyons », marque le début de la phrase. Au-dessus de « aucune »
biffé, un petit mot biffé ne se lit pas, non plus que les trois ou
quatre mots biffés écrits au-dessus de « qui ». « diverses », « pas d'évènement » sont écrits au-dessus des mots
biffés qu'ils remplacent ; « de l'ancien » au-dessus de « du ». Un
crochet ouvre « La première enseignoit » dans le texte. Des croix de
renvoi suivent « l'univers », « prétendues » et « comètes », sans objet
sur la page. « astronomie judiciaire » est écrit ainsi, après
« astrologie judiciaire » plus haut. La phrase se poursuit sur la vue
suivante.}

\page{33}\folio{16r}
n'auroient pas \struck{quitté} pris la peine de visiter la terre s'ils
n'eussent été chargés \struck{de \ill{}} d'avertir ses habitans d'un
aussi grand malheur.

\marginal{M. \ill{}}
Mais, si nous avons renoncé à ces antiques erreurs, nous conservons
cependant encore dans nos argumentations l'invincible habitude de juger
de la nature des choses par la possibilité de nous en former une idée;
en sorte qu'une proposition est affirmée ou niée suivant que nous
pouvons ou ne pouvons pas concevoir son existence. Ainsi nous disons
hardiment que la matière est divisible à l'infini; parcequ'il nous est
facile de continuer l'opération arithmétique de la division à l'infini.
nous disons qu'elle ne peut penser, parcequ'elle est divisible à
l'infini et que l'unité de nos opérations intellectuelles répugne à
l'idée de la divisibilité.

Cependant nous ne saurons toutes ces choses ni \emph{à posteriori},
puisque l'expérience ne sauroit les atteindre, ni \emph{à priori},
puisque la matière ne nous \struck{est} étant connue que par de simples
perceptions, \struck{et que par conséquent} nous ignorons complètement
son essence. On croiroit à voir notre assurance qu'à l'exemple du
géomètre nous sommes parvenus à exprimer la nature du sujet avec une
telle précision que toutes ses propriétés \struck{sont} renfermées dans
notre définition. mais combien la différence est grande! Au lieu d'une
équation absolue, qui renferme l'objet de nos recherches tout entier, en
sorte que rien de ce qui lui appartient ne puisse être étranger à cette
espèce de définition caractéristique, nous connoissons seulement quelques
propriétés relatives à nos sens. Que penser de la singulière
\note{folio 16 en haut à droite, et au-dessous, d'une encre noire et épaisse,
le chiffre romain « II » ; en haut à gauche, à l'encre, le nombre « 23 » barré d'un trait. Dans la marge de gauche, en regard de « Mais,
si nous avons renoncé », « M. » suivi d'un nom de sept ou huit lettres,
commençant par un D, souligné d'un trait qui revient vers le texte ; le
nom n'est pas lu. Un crochet ouvre l'alinéa. « est » est écrit au-dessus
de « ou affirmée », « étant » au-dessus de « est » biffé ; « à posteriori » et « à priori » sont
soulignés. Au mot biffé qui suit « chargés » sont superposées des lettres
qui ne se lisent pas. La phrase se poursuit sur la vue suivante.}

\page{34}
assurance \struck{de nos décisions} avec laquelle, lorsque nous
\uncertain{avons} à balancer \struck{sur la} les probabilités \struck{de
la formation plus ou moins grande de l'affirmation ou de la négation}
dans des questions qui ont si peu de prise, \struck{nous disons, et que
pourtant} nous n'hésitons pourtant pas à dire: il est évident, il est
absurde, il faut être de mauvaise foi pour ne pas convenir \&c.
Avouons le. La philosophie a \struck{sans doute} fait \struck{de grands}
des progrès réels mais \struck{pourtant} elle doit encore subir de
grands changemens si elle peut espérer d'arriver à l'exactitude.

Nous avons déjà remarqué qu'il existe en nous un sentiment profond
d'unité, d'ordre et de proportion qui sert de guide à tous nos jugemens.

Dans les choses morales, nous en tirons la règle du bien. Dans les
choses intellectuelles, nous y puisons la connoissance du vrai. Dans les
choses de pur agrément, nous y trouvons le caractère du beau.

Il nous est \struck{absolument impossible} difficile de savoir si les
conditions qui sont imposées à notre approbation en toutes choses, sont
le résultat \struck{nécessaire} immédiat des lois de l'être, ou si elles
dérivent seulement d'un rapport entre toute autre réalité et celle de
notre existence.

Plusieurs philosophes paroissent s'être proposé plus ou moins
directement les questions que ce doute pourroit faire naître; les uns
ont cherché dans les causes occasionnelles de nos sensations des
qualités qui leur fussent correspondantes. D'autres ont prétendu nier
l'existence des objets qui nous sont extérieurs.
\note{verso du folio 16 ; en haut à gauche, à l'encre, le nombre « 24 » barré d'un trait. Les onze premières lignes, jusqu'à « l'exactitude », sont
d'une encre plus pâle que la suite, les corrections d'une encre noire.
« avec laquelle », « les », « des », « pourtant », « Avouons le »,
« réels », « difficile » et « immédiat » sont écrits au-dessus des
lignes ; au-dessus de « lorsque », un petit mot barré suivi d'une
virgule ne se lit pas, et « avons » est surchargé. Dans
« l'exactitude », l'article est écrit sur « la ». Un crochet ouvre dans
la marge l'alinéa « Nous avons déjà remarqué ». Un trait suit « nos
jugemens ».}

\page{35}\folio{17r}
De nos jours, Kant \struck{a fort bien établi} a discuté une question de
ce genre. \struck{eux} \struck{suivant} \uncertain{sa} remarque expresse
\struck{que} les argumens les plus concluans peuvent être attribués ou
à des rapports nécessaires, ou aux formes de notre entendement: en
sorte qu'à cet égard, toute décision rationnelle \struck{nous} parfaite
\struck{être} interdite quant au raisonnement \emph{à priori}.

On ne sauroit nier en effet la légitimité de ce doute philosophique: car
\struck{il} ce doute est fondé sur l'impossibilité de comparer aucun
autre jugement à celui de l'homme.

Cependant l'opinion qui attribueroit à l'être considéré en lui même
l'unité, l'ordre et les proportions que nous cherchons dans tous les
objets \struck{de notre attention} qui fixent notre attention auroit en
sa faveur certaines inductions qu'il n'est peut-être pas inutile de
développer. Nous allons tâcher d'exposer clairement la nature de ces
inductions; il sera facile \struck{est} d'apprécier ensuite quel degré
de confiance il convient de leur accorder.

Voici quelques observations préliminaires.

Notre logique se compose \struck{d'un certain nombre} de règles dictées
par la raison universelle. Ces règles ne seroient pas moins certaines
pour nous lorsqu'on \struck{admettroit} voudroit qu'elles enseignassent
seulement à former et à reconnoître les jugemens que tout homme de bon
sens ne sauroit contester. Si nous adoptons, pour un instant
l'hypothèse de l'entier isolement de la raison, c'est-à-dire si nous
supposons qu'aucun objet extérieur à l'esprit de l'homme ne soit venu à
sa connoissance,
\note{folio 17 en haut à droite ; en haut à gauche, à l'encre, le nombre « 25 » barré d'un trait. Le début est très repris et sa syntaxe
n'est pas arrêtée. « a discuté » est écrit au-dessus de « a fort bien
établi » biffé. Au-dessus de « eux suivant » biffés, trois ou quatre mots
biffés, dont le premier est « il », et « à » ; dans la marge de droite,
deux petits mots, lus « porte par » et « ce que », sont ramenés par un
trait après « une question » et « expresse » ; « que » est biffé en fin
de ligne. La lecture « sa remarque expresse » est donc incertaine. Un
petit mot est biffé devant « argumens ». « parfaite » est écrit
au-dessus de « nous » biffé ; « ce doute » au-dessus de « il » biffé ;
« voudroit » au-dessus de « admettroit » biffé ; « enseignassent » est
écrit sur une première terminaison. « quant au raisonnement à priori »
est ajouté sous la ligne, et « à priori » souligné.}

\page{36}
\struck{et que}, livré uniquement à ses propres pensées et à \struck{celles}
qu'il \struck{qu'il doit aux} sociétés humaines, il ait cherché à
rassembler dans un corps de doctrine les vérités de son être, celles qui
naissent de ses rapports sociaux, de ses affections et de ses devoirs;
nous y trouverions \struck{que} les idées du bien, du vrai et du beau
\struck{seront telles qu'elles nous} qui nous sont actuellement connues.
Notre morale, notre logique et nos règles du goût, ne seroient
\struck{\ill{}} pas changées; car \struck{alors} même, l'art \struck{des
reproduire les impressions reçu} des récits animés, de la peinture des
passions, \struck{substitueroient encore} l'invention d'une action
poétique offriroient \struck{encore} des sujets à l'art d'embellir
\struck{et à celui} de plaire et la littérature appauvrie ne seroit
pourtant pas anéantie.

Dans cette position hypothétique, la question de savoir si les rapports
entre les différentes parties d'un sujet sont nécessaires en eux mêmes
ou s'ils nous semblent tels uniquement en vertu de nos formes
intellectuelles ne se seroit \struck{sans doute} pas présentée à
l'esprit des philosophes. Peut-être même eussent-ils été dans
l'impossibilité d'en comprendre le sens, uniquement environnés des
choses humaines; et comment en effet \struck{auroient}-ils songé à la
notion abstraite de l'être? \struck{puisqu'un} lorsqu'un seul mode
d'existence \struck{de nos} leur eût été connu? \struck{et} Leur
\struck{forme} logique \struck{eût été} \uncertain{pouvoit être} la
nôtre; \struck{mais} \struck{les conclusions} mais leurs opinions
dogmatiques eussent été \struck{pourtant} fort différentes de celles qui
ont vieilli parmi nous.
\note{verso du folio 17 ; en haut à gauche, à l'encre, le nombre « 26 » barré d'un trait.
Les deux premiers mots, « et que », sont biffés d'un trait léger ;
« uniquement » est écrit au-dessus de « livré ». La deuxième ligne
commence par « qu'il », répété, et la syntaxe de « et à celles qu'il qu'il
doit aux sociétés humaines » n'est pas arrêtée : le texte donne ici ce
que les ratures laissent. « y » est écrit au-dessus de « trouverions »,
écrit sur « trouverons ». Au-dessus de « Notre », « ne seroient » suivi
d'un mot biffé qui ne se lit pas, et « pas changées » ; « les » est écrit
au-dessus de « les » biffé devant « récits » ; « de plaire » au-dessus de
« et à celui » biffés. Un « 2 » est porté dans la marge devant « la
peinture des passions » et un « 1 » devant « des sujets », indiquant
peut-être un ordre à rétablir. « sujet » est écrit au-dessus de « d'un ».
L'addition « et comment en effet », entre les lignes, est suivie d'une
croix ; « lorsqu'un » est écrit au-dessus de « puisqu'un » biffé, et un
long trait ramène « un seul mode d'existence\ldots{} connu? » après « la
notion abstraite de l'être? ». Au-dessus de « eût été » biffés, deux mots
dont le premier est surchargé, lus avec doute « pouvoit être » ; « mais »
est écrit au-dessus de « les conclusions » biffés, et « pourtant » biffé
au-dessus de « fort ».}

\page{37}\folio{18r}
Arrêtons \struck{nous un instant} un instant notre attention sur l'objet
du doute philosophique, et tâchons d'en \struck{faire} bien définir la
nature.

La question qui a été proposée par Kant tend à sapper dans ses fondemens
la réalité absolue de toutes les certitudes que nous pouvons obtenir;
elle réduit à n'être que des vérités relatives, celles mêmes dont nous
possédons les plus claires démonstrations. Le doute que ce philosophe
célèbre attaqueroit principalement ce que nous avons admis concernant
les attributs de l'être. Ainsi le type intérieur qui nous sert à
distinguer le bien, le vrai et le beau seroit bien en effet celui qui
convient à notre manière de sentir, mais n'auroit en dehors de nous
aucune réalité dont nous pussions obtenir l'assurance.

L'auteur, après avoir dénié la preuve de l'existence de Dieu, qui se
fonde sur la nécessité de trouver une cause première à \struck{l'existence}
cette de l'univers, cherche \struck{ce qui manque} dans le sentiment ce
qui manque au raisonnement. Mais il est facile de voir que cette
concession en faveur des idées morales est purement arbitraire, et
qu'elle est \struck{uniquement} destinée à servir de sauve garde au
système des \struck{la causalité} formes intellectuelles.

Nous l'avons déjà dit, et cette proposition est fondamentale, il n'existe
qu'un seul modèle du vrai, ses copies diffèrent comme les objets qui en
reçoivent l'empreinte. Dans \struck{le bien, dans la vérité} la morale,
dans la science, dans la littérature, dans les beaux arts, nous cherchons
toujours l'unité d'existence, l'ordre et les proportions \struck{d'un même
tout. dans} entre les parties d'un même tout.
\note{folio 18 en haut à droite ; en haut à gauche, à l'encre, le nombre « 27 » barré d'un trait. Écriture plus régulière et moins
raturée que celle des pages précédentes. « phi » est récrit au-dessus de
« philosophique ». « cette », au-dessus de « l'existence » biffé, laisse
la construction « à cette de l'univers » telle quelle. Une croix de
renvoi suit « de Dieu », sans objet sur la page. Dans « des formes
intellectuelles », « des » est écrit sur « de la ». « entre » est écrit
au-dessus de « dans » biffé.}

\page{38}
Voici la question qui se présente: le modèle du vrai, ce type de l'être,
le devons nous au fait de notre existence considérée abstractivement,
c'est-à-dire, suffit-il qu'il existe un être intelligent pour qu'il
trouve en lui même les conditions sans lesquelles aucune existence n'est
possible; ou est-ce au mode particulier de notre être qu'appartiennent
les conditions qui sont pour nous le caractère du vrai?

Notre question comprend celle de Kant, \struck{celle-ci} qui pourroit
être ainsi exprimée: notre logique est-elle celle de la raison absolue,
ou convient-elle uniquement à la raison humaine?

A l'égard de ce que ce philosophe remarque touchant notre tendance
intellectuelle à chercher les causes de tout ce qui frappe notre
attention, il me paroît qu'en adoptant notre manière d'envisager les
choses, cette tendance \struck{ne} seroit l'avertissement que nous ne
voyons pas dans son entier l'objet que nous examinons. Il
\struck{se présente} s'offre à nous avec le caractère fractionnaire,
nous demandons \struck{qu'elle en} quelle en est l'unité; nous
\struck{pour} le voyons comme étant une partie; nous cherchons
\struck{à quel tout} le tout auquel cette partie appartient.

Prenons un exemple; supposons qu'au lieu d'envisager l'équation du
cercle, nous soyons frappés d'une des propriétés des sinus et cosinus;
nous pourrions bien demander pourquoi cette propriété a lieu en effet;
car alors nous n'aurions sous les yeux qu'une partie du sujet; mais, si
nous remontons jusqu'à le
\note{verso du folio 18 ; en haut à gauche, à l'encre, le nombre « 28 » barré d'un trait.
« qui » est écrit au-dessus de « celle-ci » biffé ; « s'offre » au-dessus
de « se présente » biffé ; « le » et « auquel » au-dessus de « à quel
tous », dont un petit mot est lui-même biffé : le texte donne la leçon
finale, où « tout » reste à sa place. Dans « paroît » une lettre est
surchargée d'une croix. Des croix de renvoi suivent « lui même » et
« un exemple », sans objet sur la page. La phrase se poursuit sur la
vue suivante.}

\page{39}\folio{19r}
première expression de la courbe, notre curiosité est pleinement
satisfaite; nous avons défini l'essence; nous voyons une existence
complète. Bien certainement cet être absolu et nécessaire seroit compris
de la même manière par les intelligences les plus diverses que nous
puissions imaginer.

\marginal{Pennequin}
Mais de pareils sujets sont en petit nombre. ils appartiennent aux
mathématiques pures, nos raisonnemens logiques s'appliquent au contraire,
\struck{et} à tous les sujets que nous considérons. Nous avons vu que la
question de leur certitude absolue ou relativité seroit insoluble
\emph{à priori}; que, \struck{elle seroit} si elle pouvoit être comprise
par l'homme que nous avons supposé environné uniquement des choses
humaines, il hésiteroit pas à affirmer que rien n'est plus absolu que
\struck{\ill{}} ses nécessités intellectuelles. Sans doute même il iroit
plus loin; \struck{Ainsi par exemple} et à beaucoup d'égards ses idées
seroient contraires aux nôtres.

Ainsi, par exemple, j'ai dit comment nous sommes parvenus à établir que
la matière ne pense pas. L'homme, que je suppose ne connoître autre chose
que lui même et ses semblables, n'auroit pu imaginer qu'il y eût deux
substances en lui; aucune action extérieure ne l'\struck{auroit} forçoit
à chercher des individualités invisibles et douées de volontés;
\struck{\ill{}} il \uncertain{n'eût} pas douté \struck{été pour lui un}
de l'unité de son existence; et, si, dans cette position hypothétique des
corps inertes lui eussent été présentés, il n'\struck{\ill{}} \ill{} pu
manquer de les croire doués de sentiment et de pensées; l'expérience
seule auroit \struck{pu} fini par réformer ce dogme que la matière
\struck{est douée} ou l'étendue pense et réfléchit, veut et agit.
\note{folio 19 en haut à droite ; en haut à gauche, à l'encre, un nombre barré
d'un trait, lu « 29 » avec quelque doute, le premier chiffre étant
empâté. Dans la marge de gauche, en regard de
« Mais de pareils sujets », le mot « Pennequin », d'une grande écriture,
souligné d'un trait qui revient vers le texte ; plus bas, en regard de
« ou relativité », le mot « Absolue », ici mis dans le texte après
« certitude ». « la question de » est écrit au-dessus de « leur ».
« environné » est écrit sur un premier mot ; « pas » est ajouté sous
« hésiteroit », sans « n' ». Au-dessus de « il peut été pour lui un »,
où « été pour lui un » est biffé, sont écrits « pas douté » et « de » ;
le verbe qui précède « pas douté » est surchargé et n'est pas sûr, et le
mot biffé devant « il » ne se lit pas. « iroit » est surchargé. « son »
est écrit au-dessus de « d'existence ». Après « il n' », un mot biffé et
un mot surchargé ne se lisent pas. « fini par » est écrit au-dessus de
« pu » biffé. Une croix de renvoi suit « lui même », sans objet sur la
page.}

\page{40}
Ainsi voyons-nous les enfans, qu'on a soin de préserver du contact des
objets dont ils sont environnés attribuer la volonté de les frapper au
corps dont le choc vient à les blesser. La loi du \uncertain{talion},
loi de justice innée, les porte à rendre le coup qu'ils viennent de
recevoir, et le conseil de leur nourrice, qui les y invite,
\struck{leur} est suggéré par le désir que l'enfant manifeste
naturellement d'être vengé d'une attaque qu'il regarde comme volontaire.

L'observateur qui réfléchit pense alors que l'enfant raisonne mal;
\struck{ignore encore les propriétés de la matière et se croit beaucoup
mieux instruit} tandis que ses idées dérivent immédiatement du même
sentiment d'analogie qui a porté l'homme placé dans une position
différente à des idées dogmatiques entièrement opposées.

Où trouverons nous à présent la solution de la difficulté qui nous
occupe? Les raisonnemens \emph{à priori} ne peuvent l'atteindre, puisqu'ils
sont tous formés par la raison, dont nous voulons juger la manière
d'agir: et lorsque nous \struck{\ill{}} avons recours aux preuves
extérieures, nous trouvons que, suivant la position de l'observateur,
\struck{il est conduit aux opinions} l'analogie le conduit aux opinions
les plus contraires.

Ah! si les conjectures de l'homme eussent toujours été
réalisées; si l'expérience eût sanctionné tous les systèmes qu'il a
imaginés; si, lorsqu'il avait jugé de l'impossibilité d'un fait, d'un
ordre quelconque de choses, contemporaines ou successives, l'observation
n'eût jamais démenti ses décisions théoriques, qui pourroit douter de
l'absolutisme de nos nécessités
\note{verso du folio 19 ; en haut à gauche, à l'encre, le nombre « 30 » barré d'un trait.
« Ainsi » est surchargé à l'initiale. « talion » est écrit d'une plume
chargée et sa lecture n'est pas sûre ; « choc » est écrit sur « chose ».
« raisonne mal » est écrit au-dessus de « ignore encore », qui ouvre les
deux lignes biffées ; « ses » est écrit sur « les ». « à présent » est
écrit sur un premier mot ; « à priori » est souligné. Le mot biffé
entre « nous » et « avons » ne se lit pas sûrement. L'initiale de « Ah! » est écrite sur un premier mot biffé. « avait » est écrit ainsi, et non « avoit ». Une croix de renvoi suit « blesser », sans
objet sur la page. La phrase se poursuit sur la vue suivante, hors de
ce lot.}

\end{document}
